\chapter{Характеристика исходных данных курсового проекта}
\label{cha:ish_dannie}

\section{Исходные данные для проектирования}

Электрическая сеть сооружается в Костромской области на железобетонных опорах.

\begin{figure}[h]
	\centering
	\includegraphics[width=0.7\textwidth]{inc/svg/scheme}
	\caption{Схема расположения пунктов}
	\label{fig:scheme}
\end{figure}

Питание района электроэнергией будет осуществляться от шин 220 кВ подстанции "К" работающей в составе электроэнергетической системы. Источник питания в режиме наибольших нагрузок обеспечивает полную выдачу необходимой для потребителей активной мощности, а также 87 Мвар реактивной мощности. На шинах источника питания района в режиме наибольших нагрузок обеспечивается напряжение, равно 110 \%, а в режиме наименьших нагрузок 100 \% от номинального.

Для всех пунктов наименьшая нагрузка принимается 37 \% от наибольшей; число часов использования наибольших нагрузок составляет \(4350\; \frac{\textup{ч}}{\textup{год}}\).

Вторичное напряжение на всех сооружаемых подстанциях 10 кВ.

\section{Исходные данные по климатическим условиям и нагрузкам в пунктах потребления}

Электрическая сеть сооружается в Костромской области на железобетонных опорах.

Среднеянварская температура: \(-11,8 ^oC\)

Среднегодовая температура: \(2,7\; ^oC\)

Среднеиюльская температура: \(17,6\; ^oC\)

Район по ветру: I

Район по гололёду: I

%%% Local Variables:
%%% mode: latex
%%% TeX-master: "rpz"
%%% End: